\documentclass[11pt]{article}
\usepackage[a4paper,margin=25mm]{geometry}
\usepackage{amsmath,amssymb,amsthm,mathtools}
\usepackage{graphicx}
\usepackage[hidelinks]{hyperref}
\usepackage{microtype}

\newtheorem{theorem}{Theorem}
\newtheorem{lemma}[theorem]{Lemma}
\newtheorem{corollary}[theorem]{Corollary}
\newcommand{\E}{\mathcal E}
\newcommand{\Sca}{\mathcal{SC}}
\newcommand{\Proj}{\mathcal P}
\newcommand{\B}{\mathcal B}
\newcommand{\Ran}{\operatorname{Ran}}
\newcommand{\Tr}{\operatorname{Tr}}
\newcommand{\spec}{\operatorname{spec}}
\newcommand{\coex}{\mathord{\sim}}
\newcommand{\perpco}{\mathord{\perp}}

\title{All-dimensional coexistency criteria for Hilbert-space effects}
\author{Candidate full solution packet for arXiv:2008.08559}
\date{17 August 2026}

\begin{document}
\maketitle

\begin{center}
\fbox{\parbox{0.91\linewidth}{\textbf{Status.} Candidate full solution, likely valid, subject to expert review.  Questions 6.2 and 6.3 are answered affirmatively for arbitrary Hilbert spaces.  Question 6.4 receives an exact dimension-free support-functional characterization.}}
\end{center}

\section{Source and the three questions}

The source is G.~P.~Geh\'er and P.~\v Semrl, \emph{Coexistency on Hilbert space effect algebras and a characterisation of its symmetry transformations}, arXiv:2008.08559v2, later published in \emph{Communications in Mathematical Physics} \textbf{379} (2020), 1077--1112 \cite{GS}.  The questions occur in Section~6, page~27 of the locally rendered source PDF.

\begin{center}
\includegraphics[width=\linewidth]{figures/open_problem_crop.png}
\end{center}

Let $H$ be a complex Hilbert space with $\dim H\ge2$, and
\[
  \E(H)=\{A\in\B(H):A=A^*,\ 0\le A\le I\}.
\]
Effects $A,C$ are \emph{coexistent}, written $A\coex C$, when there are effects $M,N,G$ with $A=M+G$, $C=N+G$, and $M+N+G\le I$.  Put
\[
 A^{\coex}=\{C\in\E(H):C\coex A\},\qquad A^\perp=I-A.
\]
For a rank-one projection $P$, its strength in $A$ is
\[
 \Lambda(A,P)=\max\{t\ge0:tP\le A\},
\]
and the source proves
\begin{equation}\label{eq:threshold}
 tP\coex A\quad\Longleftrightarrow\quad
 t\le F_A(P):=\Lambda(A,P)+\Lambda(I-A,P).
\end{equation}
The three source questions ask whether the finite-dimensional classification from equality of $F_A$ extends to infinite dimensions, whether $0,1\in\spec(A)$ characterizes minimal coexistency neighborhoods in infinite dimensions, and how arbitrary inclusion $A^{\coex}\subseteq B^{\coex}$ can be characterized.

\section{Proof intuition}

There are three independent mechanisms.

First, away from endpoint eigenspaces, strengths are reciprocals of inverse quadratic forms:
\[
 F_A(P_x)=\frac1{\|A^{-1/2}x\|^2}
            +\frac1{\|(I-A)^{-1/2}x\|^2}
 \quad(\|x\|=1).
\]
On any finite-dimensional real subspace of dimension at least three, these are reciprocals of positive-definite irreducible quadratic polynomials.  Unique factorization recovers the unordered pair of inverse forms.  A common form core patches the local choices globally, while endpoint eigenspaces are detected because their rank-one projections have strength sum one.

Second, inclusion $B^{\coex}\subseteq A^{\coex}$ implies that $A$ and $B$ commute.  Localizing vectors near two points of their joint spectrum converts inclusion into a two-coordinate strength inequality.  Points over spectral values $0$ and $1$ force opposite endpoint values for $B$; comparison of either endpoint with an arbitrary joint spectral point forces both $b\ge a$ and $b\le a$.

Third, $A^{\coex}$ is exactly the Minkowski sum $[0,A]+[0,I-A]$.  Its support function against a finite-rank self-adjoint test operator has a closed positive-part trace formula.  Hahn--Banach separation therefore turns arbitrary coexistency inclusion into an explicit family of scalar inequalities.

\section{Two elementary lemmas}

For an effect $T$, write
\[
 \mathcal D_T=\Ran(T^{1/2}),\qquad
 q_T(x)=\|T^{-1/2}x\|^2\quad(x\in\mathcal D_T),
\]
where the inverse is taken on $(\ker T)^\perp$.  The Busch--Gudder strength formula gives
\begin{equation}\label{eq:strengthform}
 \Lambda(T,P_x)=
 \begin{cases}
  \|x\|^2/q_T(x),&x\in\mathcal D_T,\ x\ne0,\\
  0,&x\notin\mathcal D_T,
 \end{cases}
\end{equation}
where $P_x$ projects onto $\mathbb Cx$ \cite{BG,GS}.

\begin{lemma}[reciprocal quadratic rigidity]\label{lem:ufd}
Let $V$ be a real vector space of finite dimension at least three, and let $q_1,q_2,r_1,r_2$ be positive-definite quadratic forms on $V$.  If $q_1,q_2$ are not proportional and
\[
 \frac1{q_1(x)}+\frac1{q_2(x)}
 =\frac1{r_1(x)}+\frac1{r_2(x)}\qquad(x\ne0),
\]
then $\{r_1,r_2\}=\{q_1,q_2\}$.

It is impossible on a nonempty open subset of $V$ to have
\[
 \frac1{q_1}+\frac1{q_2}=\frac1r
\]
with $r$ a positive-definite quadratic form.
\end{lemma}

\begin{proof}
A positive-definite quadratic form in at least three real variables is irreducible over $\mathbb R$.  Clearing denominators in the first identity gives
\[
 (q_1+q_2)r_1r_2=(r_1+r_2)q_1q_2.
\]
The irreducible forms $q_1,q_2$ are coprime, and neither divides $q_1+q_2$.  Unique factorization therefore matches each $q_i$, up to a positive scalar, with one $r_j$.  Substitution in the reciprocal identity makes both scalars one.  The single-reciprocal identity would give $r(q_1+q_2)=q_1q_2$; the same divisibility argument forces $r$ proportional to one $q_i$ and then forces $q_1,q_2$ proportional, a contradiction.  An identity on an open set is a polynomial identity after denominators are cleared.
\end{proof}

\begin{lemma}[one-variable pole rigidity]\label{lem:poles}
Let $0<c<1$ and $p,q,r>-1$.  If
\[
 \frac{c}{1+sp}+\frac{1-c}{1+sq}=\frac1{1+sr}
 \qquad(0<s<1),
\]
then $p=q=r$.
\end{lemma}

\begin{proof}
After clearing denominators, comparison of the coefficients of $s$ and $s^2$ says that the unordered pair $\{p,q\}$ equals
\[
 \{r,\,c q+(1-c)p\}.
\]
If $p\ne q$, the second displayed number lies strictly between them and cannot be either member of the pair.  Hence $p=q$, and then $r=p$.
\end{proof}

We also use the standard variational identity
\begin{equation}\label{eq:variational}
 \langle Ty,y\rangle
 =\sup_{x\in\mathcal D_T}
 \bigl(2\operatorname{Re}\langle y,x\rangle-q_T(x)\bigr)
 \qquad(T\ge0).
\end{equation}
It follows directly by putting $x=T^{1/2}z$ and completing the square; the formula remains valid when $T$ has a kernel.

\section{Question 6.2: rank-one coexistency determines the effect}

\begin{theorem}\label{thm:rankone}
For arbitrary $A,B\in\E(H)$, the following are equivalent:
\begin{enumerate}
\item $F_A(P)=F_B(P)$ for every rank-one projection $P$;
\item $A^{\coex}\cap\{tP:0<t\le1,\ P\in\Proj_1(H)\}
      =B^{\coex}\cap\{tP:0<t\le1,\ P\in\Proj_1(H)\}$;
\item either $A,B$ are scalar effects, or $B=A$, or $B=I-A$.
\end{enumerate}
Thus the answer to Question 6.2 is affirmative, with no separability assumption.
\end{theorem}

\begin{proof}
The equivalence of (1) and (2) follows from \eqref{eq:threshold}; (3) implies both immediately.  Assume (1).

\smallskip
\noindent\emph{The endpoint-free case.}
Suppose first that $\ker A=\ker(I-A)=\{0\}$.  Set
\[
 \mathcal D=\mathcal D_A\cap\mathcal D_{I-A}.
\]
This is a dense complex-linear subspace, since it contains the ranges of all spectral cuts $E_A([1/n,1-1/n])$.  If $A$ is non-scalar, the forms $q_A$ and $q_{I-A}$ are not proportional on $\mathcal D$: proportionality, followed by closure on the spectral-cut core, would give $A^{-1}=c(I-A)^{-1}$ and hence $A=(1+c)^{-1}I$.

If $A$ is scalar, then $F_A(P)=1$ for every $P$.  Thus every rank-one projection coexists with $B$, and therefore commutes with $B$; hence $B$ is scalar.  We may consequently assume for the remainder of this case that $A$ is non-scalar.

Choose a finite-dimensional real $V\subset\mathcal D$, $\dim_{\mathbb R}V\ge3$, on which the two forms are not proportional.  For every nonzero $x\in V$, the left side of
\begin{equation}\label{eq:recipid}
 \frac1{q_A(x)}+\frac1{q_{I-A}(x)}
 =\Lambda(B,P_x)/\|x\|^2+\Lambda(I-B,P_x)/\|x\|^2
\end{equation}
is positive.  Hence
\[
 V\subseteq\mathcal D_B\cup\mathcal D_{I-B}.
\]
The intersections of the two domains with $V$ are real linear subspaces, and a vector space over an infinite field is not the union of two proper subspaces.  Thus one domain contains $V$.  If the other did not, then on the nonempty open complement of its intersection with $V$, equation \eqref{eq:recipid} would express the sum of two reciprocal nonproportional quadratic forms as one reciprocal quadratic form, contradicting Lemma~\ref{lem:ufd}.  Consequently $V$ lies in both domains, and Lemma~\ref{lem:ufd} gives
\begin{equation}\label{eq:localpair}
 \{q_B|_V,q_{I-B}|_V\}=\{q_A|_V,q_{I-A}|_V\}.
\end{equation}

Fix one such base space $V_0$.  For any $x\in\mathcal D$, apply the same argument to the real span of $V_0$ and $x$.  Because the two forms remain nonproportional on the enlarged space, \eqref{eq:localpair} holds there.  Its orientation must agree with the one on $V_0$.  We therefore have, on all of $\mathcal D$, either
\begin{equation}\label{eq:directforms}
 q_B=q_A,\qquad q_{I-B}=q_{I-A},
\end{equation}
or the two equalities with $A$ and $I-A$ exchanged.

The spectral cuts $E_A([1/n,1-1/n])$ show that $\mathcal D$ is a form core for both inverse forms of $A$.  In the direct orientation, \eqref{eq:variational} and \eqref{eq:directforms} yield
\[
 \langle By,y\rangle\ge\langle Ay,y\rangle,\qquad
 \langle(I-B)y,y\rangle\ge\langle(I-A)y,y\rangle
\]
for every $y\in H$: the supremum for $B$ is over a domain containing $\mathcal D$, while the supremum over the core $\mathcal D$ equals the value for $A$.  Hence $B\ge A$ and $B\le A$, so $B=A$.  The exchanged orientation similarly gives $B=I-A$.

\smallskip
\noindent\emph{Endpoint eigenspaces.}
Let
\[
 H_0=\ker A,\qquad H_1=\ker(I-A),\qquad
 H_g=(H_0\oplus H_1)^\perp,
\]
and write $A_g=A|_{H_g}$.  For every unit $x\in H_0\cup H_1$, $F_A(P_x)=1$.  Therefore $P_x$ coexists with $B$ and, since $P_x$ is a projection, it commutes with $B$ \cite[Lemma~2.1(b)]{GS}.  It follows that $B$ reduces $H_0,H_1,H_g$ and is scalar on each nonzero endpoint space:
\[
 B=c_0I_{H_0}\oplus B_g\oplus c_1I_{H_1}.
\]
On $H_g$, the already proved endpoint-free case says $B_g=A_g$ or $I-A_g$ whenever $A_g$ is non-scalar.  If $A_g=eI$ is scalar, equality of the restricted strength sums makes every rank-one projection on $H_g$ commute with $B_g$, so $B_g=dI$ is scalar.

Assume first $H_0\ne\{0\}$ and $H_g\ne\{0\}$, and choose unit vectors $x\in H_0$ and
\[
 y\in\mathcal D_{A_g}\cap\mathcal D_{I-A_g}.
\]
For $z_s=\sqrt{1-s}\,x+\sqrt{s}\,y$, $0<s<1$, formula \eqref{eq:strengthform} gives
\begin{equation}\label{eq:endmix}
 F_A(P_{z_s})=
 \frac1{1+s(q_{I-A_g}(y)-1)}.
\end{equation}
If $A_g$ is non-scalar and $B_g=A_g$, then $0<c_0<1$ would turn the left side for $B$ into
\[
 \frac{c_0}{1+s(c_0q_{A_g}(y)-1)}
 +\frac{1-c_0}{1+s((1-c_0)q_{I-A_g}(y)-1)}.
\]
Lemma~\ref{lem:poles}, compared with \eqref{eq:endmix}, would imply
$(1-c_0)q_{I-A_g}(y)=q_{I-A_g}(y)$, impossible.  Thus $c_0\in\{0,1\}$.  If $c_0=1$, equation \eqref{eq:endmix} forces $q_{A_g}=q_{I-A_g}$ on a common core, hence $A_g=I/2$, again impossible.  Therefore $c_0=0$.  The alternative $B_g=I-A_g$ gives symmetrically $c_0=1$.

If $A_g=eI$ and $B_g=dI$, restriction to $\mathbb Cx\oplus\mathbb Cy$ is the two-dimensional diagonal situation already proved in Corollary~2.11 of the source: equality of all rank-one strength sums gives $(c_0,d)=(0,e)$ or $(1,1-e)$.  Thus the same global alternatives hold.

If both $H_0$ and $H_1$ are nonzero, their two-dimensional compression forces $(c_0,c_1)=(0,1)$ or $(1,0)$, consistently with the orientation on $H_g$.  If $H_g=0$, this argument is exactly the projection case.  If $H_0=0$ but $H_1\ne0$, apply the preceding argument to $I-A$, noting that $F_{I-A}=F_A$.  The only remaining cases are scalar.  This proves (3).
\end{proof}

\section{Question 6.3: minimal coexistency neighborhoods}

\begin{lemma}[joint spectral transfer]\label{lem:transfer}
Let $A,B$ be commuting effects and suppose $F_B(P)\le F_A(P)$ for every rank-one projection $P$.  Let $X$ be the Gelfand spectrum of $C^*(I,A,B)$ and write $a,b\in C(X)$ for the Gelfand transforms.  Then for all $s,t\in X$ and $0<\theta<\pi/2$,
\begin{equation}\label{eq:twopoint}
 T_{b(s),b(t)}(\theta)\le T_{a(s),a(t)}(\theta),
\end{equation}
where
\[
 T_{u,v}(\theta)=
 \left(\frac{\cos^2\theta}{u}+\frac{\sin^2\theta}{v}\right)^{-1}
 +\left(\frac{\cos^2\theta}{1-u}+\frac{\sin^2\theta}{1-v}\right)^{-1},
\]
with a reciprocal equal to zero when its denominator is infinite.
\end{lemma}

\begin{proof}
For distinct $s,t$, choose shrinking disjoint joint spectral neighborhoods on which both $a$ and $b$ differ from their values at the centers by at most $\varepsilon$.  Their joint spectral projections in $C^*(A,B)''$ are nonzero: otherwise a nonzero continuous function supported in such a neighborhood would be represented by zero, contradicting faithfulness.  Choose unit vectors in the two orthogonal reducing ranges and form their $\theta$-mixture.  On each range, $A$ and $B$ are squeezed in operator order between the corresponding scalar values plus or minus $\varepsilon$.  Apply \eqref{eq:strengthform} separately to the effect and its complement.  Letting $\varepsilon\downarrow0$ gives \eqref{eq:twopoint}.  The case $s=t$ follows by continuity or is immediate.
\end{proof}

\begin{theorem}\label{thm:minimal}
For every effect $A$ on an arbitrary Hilbert space $H$, the following are equivalent:
\begin{enumerate}
\item $0,1\in\spec(A)$;
\item there is no $B\in\E(H)$ such that $B^{\coex}\subsetneq A^{\coex}$.
\end{enumerate}
Thus Question 6.3 has an affirmative answer, without separability.
\end{theorem}

\begin{proof}
Assume $0,1\in\spec(A)$ and $B^{\coex}\subseteq A^{\coex}$.  Lemma~2.12(a) of the source, with the roles there chosen as $B,A$, gives $A\in B''$; in particular $A$ and $B$ commute \cite{GS}.  By \eqref{eq:threshold}, inclusion gives $F_B\le F_A$.

Use the notation of Lemma~\ref{lem:transfer}.  Since $a(X)=\spec(A)$, choose $s_0,s_1\in X$ with $a(s_0)=0$ and $a(s_1)=1$.  For every interior $\theta$,
\[
 T_{0,1}(\theta)=0.
\]
Equation \eqref{eq:twopoint} and nonnegativity imply
\[
 \{b(s_0),b(s_1)\}=\{0,1\}.
\]
Suppose first $b(s_0)=0$ and $b(s_1)=1$.  For arbitrary $s\in X$, comparison with $s_0$ gives
\[
 T_{0,b(s)}(\theta)\le T_{0,a(s)}(\theta).
\]
For interior $\theta$,
\[
 T_{0,u}(\theta)=
 \left(\cos^2\theta+\frac{\sin^2\theta}{1-u}\right)^{-1},
\]
which is decreasing in $u$; hence $b(s)\ge a(s)$.  Comparison with $s_1$ gives
\[
 T_{1,b(s)}(\theta)\le T_{1,a(s)}(\theta),\qquad
 T_{1,u}(\theta)=
 \left(\cos^2\theta+\frac{\sin^2\theta}{u}\right)^{-1},
\]
and this is increasing in $u$, so $b(s)\le a(s)$.  Thus $b=a$ on $X$ and $B=A$.  The reversed endpoint orientation gives $b=1-a$ and $B=I-A$.  In either case the inclusion is equality.

Conversely, if an endpoint is missing and $A$ is non-scalar, the scaling argument from the proof of Lemma~6.1 in the source applies verbatim in every dimension.  For example, if $\spec(A)\subset[0,1-\varepsilon]$, then a suitable nontrivial rescaling $C=(1-\varepsilon)^{-1}A$ satisfies $C^{\coex}\subsetneq A^{\coex}$ by convexity and Theorem~2.9 of the source.  The case with the other endpoint is obtained from $I-A$.  If $A$ is scalar, $A^{\coex}=\E(H)$ and any nontrivial projection has a proper coexistency neighborhood.  This proves the contrapositive.
\end{proof}

\section{Question 6.4: an exact support-functional characterization}

For a finite-rank self-adjoint operator $T$, define
\begin{equation}\label{eq:support}
 h_A(T)=
 \Tr\bigl((A^{1/2}TA^{1/2})_+\bigr)
 +\Tr\bigl(((I-A)^{1/2}T(I-A)^{1/2})_+\bigr),
\end{equation}
where $S_+=(|S|+S)/2$.

\begin{theorem}\label{thm:inclusion}
For arbitrary effects $A,B$ on any Hilbert space,
\begin{equation}\label{eq:characterization}
 A^{\coex}\subseteq B^{\coex}
 \quad\Longleftrightarrow\quad
 h_A(T)\le h_B(T)
 \quad\hbox{for every finite-rank self-adjoint }T.
\end{equation}
Equivalently, inclusion is determined by explicit positive-part trace inequalities on all finite-dimensional compressions.
\end{theorem}

\begin{proof}
The coexistency decomposition used in the source is precisely
\begin{equation}\label{eq:minkowski}
 A^{\coex}=[0,A]+[0,I-A],
\end{equation}
where $[0,A]=\{M:0\le M\le A\}$.  Indeed, $C=M+N$ with $M\le A$ and $N\le I-A$ is equivalent to $C\coex A$.

Fix finite-rank self-adjoint $T$.  By Douglas factorization, every $M\in[0,A]$ has the form $M=A^{1/2}KA^{1/2}$ with $0\le K\le I$ on the support of $A$, and conversely every such operator belongs to $[0,A]$.  Therefore
\begin{align*}
 \sup_{M\in[0,A]}\Tr(TM)
 &=\sup_{0\le K\le I}\Tr\bigl(KA^{1/2}TA^{1/2}\bigr)\\
 &=\Tr\bigl((A^{1/2}TA^{1/2})_+\bigr),
\end{align*}
the maximum being attained at the positive spectral projection.  Support functions add under Minkowski sums, so \eqref{eq:minkowski} gives
\[
 \sup_{C\in A^{\coex}}\Tr(TC)=h_A(T).
\]

Inclusion implies the inequalities in \eqref{eq:characterization}.  Conversely, $B^{\coex}$ is convex and WOT compact \cite[Lemma~2.6]{GS}.  If some $C\in A^{\coex}$ were outside $B^{\coex}$, real WOT Hahn--Banach separation would provide a WOT-continuous real functional strictly separating $C$ from $B^{\coex}$.  Every such functional on the self-adjoint part of $\B(H)$ is $X\mapsto\Tr(TX)$ for a finite-rank self-adjoint $T$.  This would give $h_A(T)>h_B(T)$, contradicting the assumed inequalities.  Hence inclusion holds.
\end{proof}

The criterion \eqref{eq:characterization} is a complete answer in convex-dual form.  It is not claimed to be a simpler pointwise functional-calculus classification of all pairs $(A,B)$; rather, it replaces quantified coexistency decompositions by explicit scalar spectral expressions and is valid without dimension or separability restrictions.

\section{Verification and novelty status}

\textbf{Internal checks.}  The proof was checked against the following edge cases: scalar effects; nontrivial projections; one-dimensional generic spectral summands; effects with continuous spectrum touching both endpoints; and nonseparable spaces.  The file \texttt{code/check\_identities.py} symbolically verifies the pole identity used in Lemma~\ref{lem:poles} and numerically checks the support-function maximizer and the direct/complement strength identities on random finite matrices.  These checks are diagnostics, not substitutes for the proof.

\textbf{Novelty search.}  The bounded search covered the run indexes, the locally parsed arXiv corpus, the exact title and arXiv id, the source authors, and phrases combining ``coexistence/coexistency'', ``rank-one effects'', ``infinite dimensional'', and ``minimal coexistency neighborhood''.  Later local papers citing Geh\'er--\v Semrl were inspected at the citation context; none stated an answer to Questions 6.2--6.4.  A fresh external-network search was unavailable in the current environment.  Accordingly the proof verdict is \emph{likely valid}, while novelty confidence is \emph{moderate} pending a human literature refresh.

\textbf{Recommended expert review.}  Check especially (i) the form-domain union and orientation-patching argument in Theorem~\ref{thm:rankone}, (ii) the passage from joint spectral neighborhoods to \eqref{eq:twopoint}, and (iii) the WOT separation claim in Theorem~\ref{thm:inclusion}.  If these pass, the packet resolves both infinite-dimensional questions and supplies an exact answer to the general inclusion question.

\begin{thebibliography}{9}
\bibitem{BG}
P.~Busch and S.~Gudder,
Effects as functions on projective Hilbert space,
\emph{Letters in Mathematical Physics} \textbf{47} (1999), 329--337.

\bibitem{GS}
G.~P.~Geh\'er and P.~\v Semrl,
Coexistency on Hilbert space effect algebras and a characterisation of its symmetry transformations,
arXiv:2008.08559v2; \emph{Communications in Mathematical Physics} \textbf{379} (2020), 1077--1112.
\end{thebibliography}

\end{document}
